% =====================================================================
%  System Reference: Skills
%  Auto-generated from skills/*/SKILL.md and proto/*/v1/*.proto.
%  Regenerate when those change.
%
%  Composite/atomic skills are listed once. Every leaf gRPC method is
%  listed as its own skill (one RPC = one skill), with the request and
%  response fields drawn from its proto contract.
% =====================================================================
%
% Self-contained --- no extra packages required.
%
% \skillcard{name}{inputs}{outputs}{description}


The system is built from a single kind of component, the \emph{skill}.
Every skill is declared by the same metadata schema --- a name, a
description, and its typed inputs and outputs. Higher-level
\emph{composite} skills are materializable subgraphs authored by
composing other skills; \emph{atomic} skills are single classes; and
\emph{primitive} skills are individual gRPC methods, each wrapping one
model invocation or geometric routine. All skills share a common type
system (\texttt{Se3Pose}, \texttt{OrientedBoundingBox},
\texttt{PointCloud}, \texttt{Mask}, \texttt{Image}, \texttt{Trajectory},
\texttt{JointState}, $\dots$), which is what lets them compose freely. In
the input lists below, defaults are given in parentheses and the
universal caller-context field is omitted.

% =====================================================================
\subsection{Composite and atomic skills}




\paragraph{Perception.}

\skillcard{perception\_single}
{\texttt{object\_name} (literal)}
{\texttt{<name>\_obb}: OrientedBoundingBox, \texttt{<name>\_mask}: Mask, \texttt{<name>\_cloud}: PointCloud}
{Fast single-path detection: broad detect + a VLM crop tournament +
segmentation + depth-to-3D fusion, closed by an OBB fit. Best for
uncluttered scenes with visually distinct targets.}

\skillcard{perception\_multi}
{\texttt{object\_name} (literal)}
{\texttt{<name>\_obb}: OrientedBoundingBox, \texttt{<name>\_mask}: Mask}
{Three detectors run on the same observation and a VLM disambiguator picks
the best mask on disagreement. Most robust, but slower; single-shot only.}

\skillcard{perception\_any}
{\texttt{object\_name} (literal)}
{\texttt{<name>\_obb}: OrientedBoundingBox, \texttt{<name>\_mask}: Mask, \texttt{<name>\_cloud}: PointCloud}
{Lightweight one-shot set-of-marks pick. Returns ``not found'' cleanly
when the VLM replies ``none'', the right choice for clean-all-items loops.}

\skillcard{perception\_subpart}
{\texttt{parent\_prompt}, \texttt{subpart\_prompt} (literals)}
{\texttt{<name>\_obb}: OrientedBoundingBox, \texttt{<name>\_mask}: Mask, \texttt{<name>\_cloud}: PointCloud}
{Hierarchical perception: detect the parent, crop to its box, segment the
named subpart inside the crop, uncrop and fuse depth. For affordances that
occupy few full-image pixels (pan handle, drawer pull, mug rim).}

\skillcard{track\_object}
{text prompt; observation stream}
{a stream of tracked-mask snapshots}
{Long-running tracker: seeds on the first frame via text prompt, then polls
the observation stream and advances the tracker until the workflow signals
termination, freeing tracker state on exit.}

\paragraph{Grasping.}

\skillcard{grasp\_curobo\_obb}
{\texttt{target\_obb}: OrientedBoundingBox, \texttt{target\_mask}: Mask}
{\texttt{ee\_pose\_at\_grasp}: Se3Pose, \texttt{grasp\_pose}: Se3Pose}
{Collision-aware grasping: plans a collision-free trajectory to one of
several top-down OBB grasp candidates. The default when the OBB centroid
is graspable.}

\skillcard{grasp\_moe}
{\texttt{target\_obb}: OrientedBoundingBox, \texttt{target\_mask}: Mask, \texttt{target\_cloud}: PointCloud}
{\texttt{ee\_pose\_at\_grasp}: Se3Pose}
{Mixture-of-experts grasping: a diffusion sampler and a dense OBB sweep
feed one discriminator; the top-K union is planned. Use when the OBB
centroid is not the graspable point (bowl rim, handle, off-center).}

\skillcard{grasp\_short\_axis}
{\texttt{target\_obb}: OrientedBoundingBox, \texttt{target\_mask}: Mask, \texttt{base\_obb}: OrientedBoundingBox (optional)}
{\texttt{ee\_pose\_at\_grasp}: Se3Pose, \texttt{grasp\_pose}: Se3Pose}
{Deterministic geometry-locked grasp: the finger axis snaps to the OBB's
shorter horizontal axis so the jaws close across the narrow dimension of
an elongated target. An optional base OBB slides the grasp along a handle.}

\skillcard{grasp\_direct\_ik}
{\texttt{target\_obb}: OrientedBoundingBox}
{(grasp completion only)}
{Planner-free grasping: pre-rotate to the grasp orientation above the
target, descend straight down, close. For uncluttered scenes or platforms
without a motion planner.}

\paragraph{Transport, policy, and bimanual.}

\skillcard{transport\_to\_drop}
{\texttt{container\_obb}, \texttt{container\_mask}, \texttt{target\_obb}, \texttt{target\_mask}, \texttt{ee\_pose\_at\_grasp}}
{(placement completion only)}
{Moves a held object above a destination container and releases it. An
optional VLM-grounded step localizes a natural-language sub-region (``the
left compartment'') before the drop pose is computed.}

\skillcard{run\_policy}
{observation stream; \texttt{max\_windows}}
{(policy rollout completion)}
{Runs a learned VLA policy in closed loop, reading the observation stream
each window, until a VLM termination check fires or \texttt{max\_windows}
is reached.}

\skillcard{bimanual\_crate\_lift}
{(none --- waypoints are baked parameters)}
{\texttt{done}: bool, \texttt{final\_left\_world\_pos}, \texttt{final\_right\_world\_pos}}
{Two-arm crate lift: both arms side-grasp the crate's handle bars and raise
it in lock-step, in a fixed phase order
(open\,$\to$\,approach\,$\to$\,insert\,$\to$\,close\,$\to$\,lift\,$\to$\,settle).}

% =====================================================================
\subsection{Primitive skills (one per gRPC method)}

\paragraph{Detection and segmentation.}

\skillcard{grounding\_dino.Detect}
{\texttt{image}: Image; \texttt{text\_prompt} (period-separated phrases); \texttt{box\_threshold} (0.20); \texttt{text\_threshold} (0.20)}
{\texttt{detections}: [\,box: BoundingBox2D, label, score\,]}
{Zero-shot object detection from a text prompt; returns boxes, labels, and
confidence scores.}

\skillcard{owlvit.Detect}
{\texttt{image}: Image; \texttt{text\_queries}: string[]; \texttt{score\_threshold} (0.05)}
{\texttt{detections}: [\,box: BoundingBox2D, label, score\,]}
{Open-vocabulary object detection from a list of text queries; an
alternative to grounding\_dino.}

\skillcard{sam3.SegmentText}
{\texttt{image}: Image; \texttt{text\_prompt}}
{\texttt{masks}: Mask[], \texttt{scores}: float[], \texttt{boxes}: BoundingBox2D[]}
{Segment all instances matching a text description.}

\skillcard{sam3.SegmentPoint}
{\texttt{image}: Image; \texttt{pixel\_x}, \texttt{pixel\_y}}
{\texttt{masks}: Mask[], \texttt{scores}: float[]}
{Segment the object under a single pixel coordinate.}

\skillcard{sam3.SegmentBox}
{\texttt{image}: Image; \texttt{box}: BoundingBox2D; \texttt{pixel\_x}/\texttt{pixel\_y} + \texttt{use\_point} (false, optional foreground point)}
{\texttt{masks}: Mask[], \texttt{scores}: float[]}
{Segment within a bounding box, optionally refined by a foreground point.}

\skillcard{sam2.SegmentBox}
{\texttt{image}: Image; \texttt{box}: BoundingBox2D; \texttt{max\_masks} (0 = unlimited)}
{\texttt{masks}: Mask[], \texttt{scores}: float[]}
{Box-prompted segmentation (lighter SAM2 backend).}

\skillcard{sam2.SegmentPoint}
{\texttt{image}: Image; \texttt{pixel\_x}, \texttt{pixel\_y}}
{\texttt{masks}: Mask[], \texttt{scores}: float[]}
{Point-prompted segmentation (lighter SAM2 backend).}

\skillcard{sam3\_tracker.InitTracker}
{\texttt{image}: Image (first frame); one of \texttt{text} / \texttt{box} / point (\texttt{pixel\_x},\texttt{pixel\_y},\texttt{use\_point}); \texttt{object\_name}}
{\texttt{tracker\_id}, \texttt{initial\_mask}: Mask, \texttt{initial\_box}, \texttt{score}, \texttt{object\_present}}
{Initialize a stateful tracker session seeded with one frame and a prompt.}

\skillcard{sam3\_tracker.UpdateTracker}
{\texttt{tracker\_id}; \texttt{image}: Image (new frame)}
{\texttt{mask}: Mask, \texttt{box}, \texttt{confidence}, \texttt{object\_present}}
{Advance an existing tracker session by one frame.}

\skillcard{sam3\_tracker.CloseTracker}
{\texttt{tracker\_id}}
{(empty)}
{Free a tracker session; idempotent.}

\skillcard{sam3d.GenerateMesh}
{\texttt{object\_name}; \texttt{views}: [rgb, mask, intrinsics, camera\_pose]; \texttt{obb}: OrientedBoundingBox; \texttt{view\_selection}; \texttt{seed} (42); \texttt{output\_dir}}
{\texttt{mesh\_path}, \texttt{texture\_path}, \texttt{world\_pose}: Se3Pose, \texttt{scale}: Vec3, \texttt{chosen\_view\_index}, \texttt{num\_vertices}, \texttt{num\_faces}}
{Reconstruct a triangle mesh from an RGB image + mask, scaled and posed
into world frame via the OBB; writes an \texttt{.obj} for the rehearsal sandbox.}

\paragraph{Vision--language.}

\skillcard{molmo.PointPrompt}
{\texttt{image}: Image; \texttt{object\_name}}
{\texttt{pixel\_x}, \texttt{pixel\_y}, \texttt{found}}
{Point to a named object; returns pixel coordinates.}

\skillcard{molmo.Query}
{\texttt{prompt}; \texttt{image}: Image (optional); \texttt{use\_multiview}}
{\texttt{text}}
{Free-form visual question answering (Molmo backend).}

\skillcard{molmo.QueryYesNo}
{\texttt{prompt}; \texttt{image}: Image (optional); \texttt{use\_multiview}}
{\texttt{answer}: bool}
{Yes/no visual question answering (Molmo backend).}

\skillcard{vlm.Query}
{\texttt{prompt}; \texttt{image}: Image (optional)}
{\texttt{text}}
{Free-form visual question answering (configurable VLM backend).}

\skillcard{vlm.QueryYesNo}
{\texttt{prompt}; \texttt{image}: Image (optional)}
{\texttt{answer}: bool}
{Yes/no visual question answering (configurable VLM backend).}

\paragraph{Grasp generation.}

\skillcard{graspgen.PlanFromPointCloud}
{\texttt{target\_cloud}: PointCloud; \texttt{scene\_cloud} (optional); \texttt{num\_grasps} (200); \texttt{topk\_num\_grasps} (50); \texttt{grasp\_threshold} ($-1$); \texttt{remove\_outliers} (true)}
{\texttt{grasp\_poses}: Se3Pose[], \texttt{scores}: float[], \texttt{contact\_points}: Vec3[]}
{Diffusion-based 6-DoF grasp sampling from an object-segmented cloud;
ranked, world-frame, fingertip convention.}

\skillcard{graspgen.ScoreGrasps}
{\texttt{target\_cloud}: PointCloud; \texttt{candidate\_poses}: Se3Pose[]}
{\texttt{scores}: float[] (aligned to candidates)}
{Score caller-provided grasps with GraspGen's discriminator, so
OBB-sampled candidates are comparable with diffusion samples.}

\skillcard{graspnet.PlanFromDepth}
{\texttt{depth}: DepthImage; \texttt{intrinsics}: Matrix3x3; \texttt{segmentation}: Mask; \texttt{instance\_id}; knobs (\texttt{z\_min} 0.2, \texttt{z\_max} 2.0, \texttt{max\_retries} 7, $\dots$); optional wrist camera}
{\texttt{grasp\_poses}: Se3Pose[], \texttt{scores}: float[], \texttt{contact\_points}: Vec3[]}
{Contact-GraspNet 6-DoF grasp detection from a depth image + instance mask.}

\skillcard{graspnet.PlanFromPointCloud}
{\texttt{full\_cloud}: PointCloud; \texttt{segment\_cloud}: PointCloud; \texttt{instance\_id}; same knobs}
{\texttt{grasp\_poses}: Se3Pose[], \texttt{scores}: float[], \texttt{contact\_points}: Vec3[]}
{Contact-GraspNet grasp detection from precomputed clouds.}

\skillcard{m2t2.PlanFromPointCloud}
{\texttt{scene\_cloud}: PointCloud; \texttt{target\_cloud} (optional); \texttt{num\_points} (16384); \texttt{grasp\_threshold} (0.035); \texttt{max\_grasps} (200)}
{\texttt{grasp\_poses}: Se3Pose[], \texttt{scores}: float[], \texttt{contact\_points}: Vec3[]}
{M2T2 6-DoF grasp generation from a colored scene cloud; an alternative
candidate source.}

\paragraph{Motion planning and IK.}

\skillcard{curobo.PlanToGraspPoses}
{\texttt{world\_config}: WorldConfig; \texttt{start\_joint\_position}: JointState; \texttt{grasp\_poses}: Se3Pose[]; thresholds + collision knobs}
{\texttt{success}, \texttt{trajectory}: Trajectory, \texttt{goalset\_index}}
{Plan a collision-free trajectory to one of several grasp poses (goalset).}

\skillcard{curobo.PlanWithGraspedObject}
{\texttt{world\_config}; \texttt{start\_joint\_position}; \texttt{target\_pose}: Se3Pose; \texttt{object\_name}}
{\texttt{success}, \texttt{trajectory}: Trajectory}
{Plan a collision-free trajectory while holding a named grasped object.}

\skillcard{curobo.PlanLinear}
{\texttt{start\_pose}, \texttt{end\_pose}: Se3Pose; \texttt{start\_joint\_position}; \texttt{world\_config} (optional); \texttt{num\_waypoints} (20); \texttt{hold\_vec\_weight}}
{\texttt{success}, \texttt{trajectory}: Trajectory, \texttt{failure\_reason}}
{GPU-accelerated straight-line Cartesian motion via per-waypoint IK,
optionally constrained and collision-checked.}

\skillcard{curobo.PlanDirectedLinear}
{\texttt{start\_joint\_position}; \texttt{start\_pose}, \texttt{target\_pose}; \texttt{allowed\_axes}; \texttt{endpoint\_mode}; \texttt{orientation\_mode}}
{\texttt{success}, \texttt{trajectory}: Trajectory, \texttt{failure\_reason}}
{Constrained linear motion that holds selected Cartesian axes
(project-to-target / fixed-distance / orient-in-place).}

\skillcard{curobo.PlanGraspMotion}
{\texttt{start\_joint\_position}; \texttt{grasp\_pose}: Se3Pose; approach axis+distance (0.12); lift axis+distance (0.20)}
{\texttt{success}; \texttt{approach\_trajectory}, \texttt{grasp\_trajectory}, \texttt{lift\_trajectory}: Trajectory; \texttt{failure\_reason}}
{Three-phase grasp plan (free-space approach, constrained grasp,
constrained lift) so the caller can interleave gripper commands. (CuRobo v0.8.)}

\skillcard{curobo.SolveIK}
{\texttt{target\_pose}: Se3Pose; \texttt{seed\_config} (optional); \texttt{tcp\_offset}; \texttt{num\_seeds} (32)}
{\texttt{success}, \texttt{joint\_config}: JointState}
{Pure geometric IK for a single TCP pose (no world collision).}

\skillcard{curobo.PlanToPose}
{\texttt{target\_pose}: Se3Pose; \texttt{start\_joint\_position}; \texttt{world\_config} (optional); \texttt{tcp\_offset}}
{\texttt{success}, \texttt{trajectory}: Trajectory}
{Collision-aware trajectory to a single EE target pose.}

\skillcard{curobo.BatchGraspFeasibility}
{\texttt{world\_config}; \texttt{start\_state}: JointState; \texttt{grasp\_poses}: Se3Pose[]; \texttt{approach\_offset\_m} (0.10); \texttt{ignore\_obstacle\_names}}
{\texttt{feasible}: bool[], \texttt{grasp\_ik\_ok}: bool[], \texttt{approach\_ik\_ok}: bool[], \texttt{corridor\_collision\_fraction}: float[] (all aligned to input)}
{Per-pose scene-collision feasibility for a batch of grasp candidates;
filters blocked grasps before planning.}

\skillcard{curobo.ValidateJointTrajectoryRobot}
{\texttt{world\_config}; \texttt{trajectory}: Trajectory; collision knobs}
{\texttt{success}, \texttt{failure\_reason}, \texttt{first\_collision\_waypoint}, \texttt{collision\_status\_detail}}
{Validate joint waypoints against world + self-collision.}

\skillcard{curobo.ValidateJointTrajectoryGrasped}
{\texttt{world\_config}; \texttt{trajectory}; \texttt{object\_name}; collision knobs}
{\texttt{success}, \texttt{failure\_reason}, \texttt{first\_collision\_waypoint}, \texttt{collision\_status\_detail}}
{As above, with a grasped object attached at the first waypoint.}

\skillcard{pyroki.SolveIK}
{\texttt{target\_pose}: Se3Pose; \texttt{target\_link\_name} (panda\_hand); \texttt{prev\_config} (optional); \texttt{tcp\_offset} ($0,0,-0.1$)}
{\texttt{joint\_config}: JointState, \texttt{success}}
{CPU differentiable IK for a target EE pose; the planner-free IK fallback.}

\skillcard{pyroki.PlanLinear}
{\texttt{start\_pose}, \texttt{end\_pose}: Se3Pose; \texttt{num\_waypoints} (20); \texttt{dt} (0.02)}
{\texttt{trajectory}: Trajectory, \texttt{success}}
{CPU linear Cartesian trajectory between two poses.}

\paragraph{Geometry.}

\skillcard{geometry\_svc.FilterNoise}
{\texttt{point\_cloud}: PointCloud; \texttt{eps} (0.005); \texttt{min\_samples} (10)}
{PointCloud}
{DBSCAN noise filtering of a point cloud.}

\skillcard{geometry\_svc.ComputeOBB}
{\texttt{point\_cloud}: PointCloud}
{OrientedBoundingBox}
{Fit an oriented bounding box to 3D points.}

\skillcard{geometry\_svc.FilterAndComputeOBB}
{\texttt{point\_cloud}: PointCloud; \texttt{eps} (0.005); \texttt{min\_samples} (10)}
{OrientedBoundingBox}
{FilterNoise + ComputeOBB in one round trip.}

\skillcard{geometry\_svc.TopDownGraspFromOBB}
{\texttt{obb}: OrientedBoundingBox; \texttt{z\_offset}}
{Se3Pose}
{A single top-down grasp pose from an OBB.}

\skillcard{geometry\_svc.TopDownGraspCandidates}
{\texttt{obb}: OrientedBoundingBox; \texttt{z\_offset}}
{\texttt{poses}: Se3Pose[] (primary + 90$^\circ$-rotated)}
{Candidate top-down grasps so a planner can pick a reachable one.}

\skillcard{geometry\_svc.SelectTopDownGrasp}
{\texttt{grasp\_poses}: Se3Pose[]; \texttt{gripper\_position}: Vec3}
{Se3Pose}
{Pick the most top-down (and nearest) grasp from candidates.}

\skillcard{geometry\_svc.FrontGraspFromOBB}
{\texttt{obb}: OrientedBoundingBox; \texttt{approach\_offset} (0.08); \texttt{approach\_hint}; \texttt{z\_offset}}
{\texttt{grasp\_pose}, \texttt{pre\_grasp\_pose}: Se3Pose; \texttt{approach\_direction}, \texttt{slide\_axis}: Vec3}
{Front/horizontal grasp for a handle (drawers, doors), with a slide axis.}

\skillcard{geometry\_svc.DepthToPointCloud}
{\texttt{depth}: DepthImage; \texttt{intrinsics}: Matrix3x3}
{PointCloud (camera frame)}
{Back-project a depth image to a camera-frame point cloud.}

\skillcard{geometry\_svc.MaskToWorldPoints}
{\texttt{mask}: Mask; \texttt{depth}: DepthImage; \texttt{intrinsics}: Matrix3x3; \texttt{camera\_pose}: Se3Pose}
{PointCloud (world frame)}
{Back-project a 2D mask to world-frame 3D points.}

\skillcard{geometry\_svc.PixelToWorldPoint}
{\texttt{pixel\_x}, \texttt{pixel\_y}; \texttt{depth}; \texttt{intrinsics}; \texttt{camera\_pose}}
{Vec3 (world point)}
{Back-project a single pixel to a world point.}

\skillcard{geometry\_svc.TransformPoints}
{\texttt{point\_cloud}: PointCloud; \texttt{transform}: Se3Pose}
{PointCloud}
{Apply a rigid transform to a set of points.}

\skillcard{geometry\_svc.BuildWorldConfig}
{\texttt{cameras}: CameraObservation[]; \texttt{object\_masks}: [name, mask, camera\_index]; recon params (\texttt{voxel\_size} 0.005, \texttt{mesh\_alpha} 0.03); \texttt{robot\_joint\_state} (optional); \texttt{target\_obb} (optional)}
{\texttt{config}: WorldConfig, \texttt{mesh\_names}: string[]}
{Reconstruct a named collision-mesh world from camera observations for any
planner.}

\skillcard{geometry\_svc.ComputeDropPosition}
{\texttt{container\_obb}: OrientedBoundingBox; \texttt{clearance} (0.05); \texttt{object\_z\_extent}}
{Vec3 (drop point)}
{Drop position above a container, accounting for clearance and object height.}

\skillcard{geometry\_svc.ComputeXYDistance}
{\texttt{point\_a}, \texttt{point\_b}: Vec3}
{\texttt{distance}: double}
{XY-plane Euclidean distance between two points.}

\skillcard{geometry\_svc.RotateQuatZ90}
{\texttt{quat}: Quaternion (WXYZ)}
{Quaternion}
{Rotate a quaternion 90$^\circ$ about Z (for grasp candidates).}

\paragraph{Robot and environment control.}

\skillcard{robot\_control.GetEEPose}
{\texttt{arm\_id} (0)}
{\texttt{pose}: Se3Pose}
{Current end-effector pose in world frame.}

\skillcard{robot\_control.GoToPose}
{\texttt{pose}: Se3Pose; \texttt{z\_approach}; \texttt{tcp\_offset} ($0,0,-0.1$); \texttt{tolerance} (0.01); \texttt{max\_steps} (120); \texttt{nonblocking}}
{(empty)}
{Move the EE to a target pose via IK (no collision avoidance).}

\skillcard{robot\_control.GoToPoseCartesian}
{\texttt{pose}: Se3Pose; \texttt{lin\_vel\_norm} (1.0); \texttt{ang\_vel\_norm} (2.0); \texttt{timeout\_s} (20)}
{(empty)}
{Move to a pose with smooth Cartesian interpolation.}

\skillcard{robot\_control.ExecuteJointTrajectory}
{\texttt{trajectory}: Trajectory; \texttt{subsample} (1); \texttt{tolerance} (0.01)}
{(empty)}
{Execute a pre-planned joint trajectory (from CuRobo or PyRoKI).}

\skillcard{robot\_control.MoveToJoints}
{\texttt{joint\_config}: JointState; \texttt{tolerance} (0.01); \texttt{max\_steps} (120)}
{(empty)}
{Move directly to a joint configuration.}

\skillcard{robot\_control.GoHome}
{(none)}
{(empty)}
{Move the robot to its safe home configuration.}

\skillcard{robot\_control.GoToPoseArm}
{\texttt{arm\_id}; \texttt{pose}: Se3Pose; \texttt{z\_approach}; \texttt{tcp\_offset}}
{(empty)}
{Multi-arm: move one specified arm to a pose.}

\skillcard{robot\_control.GoToPoseBoth}
{\texttt{pose\_arm0}, \texttt{pose\_arm1}: Se3Pose; \texttt{z\_approach}; \texttt{tcp\_offset}}
{(empty)}
{Multi-arm: move both arms to their poses simultaneously.}

\skillcard{robot\_control.MoveToJointsBoth}
{\texttt{joints\_arm0}, \texttt{joints\_arm1}: JointState; \texttt{tolerance} (0.02); \texttt{max\_steps} (300)}
{(empty)}
{Multi-arm: drive both arms to raw joint targets in lock-step (no IK).}

\skillcard{robot\_control.ApplyPolicyAction}
{\texttt{action}: double[] (e.g. 7-dim $[\Delta x,\Delta y,\Delta z,\Delta r_x,\Delta r_y,\Delta r_z,\text{grip}]$); \texttt{arm\_id} (0)}
{(empty)}
{Apply one low-level policy action to the env controller (VLA rollouts).}

\skillcard{gripper.Open}
{\texttt{arm\_id} (0); \texttt{settle\_steps} (40)}
{GripperState}
{Open the gripper and settle.}

\skillcard{gripper.Close}
{\texttt{arm\_id} (0); \texttt{settle\_steps} (60)}
{GripperState}
{Close the gripper and settle.}

\skillcard{gripper.GetPosition}
{\texttt{arm\_id} (0)}
{GripperState (opening width)}
{Read the current gripper opening.}

\skillcard{gripper.GetPose}
{\texttt{arm\_id} (0)}
{\texttt{pose}: Se3Pose}
{Gripper pose in world frame.}

\skillcard{observation.GetObservation}
{(none)}
{\texttt{cameras}: CameraObservation[] (RGB-D + intrinsics/extrinsics), \texttt{arm\_states}: ArmState[]}
{Full sensor observation from all cameras and arms.}

\skillcard{observation.GetCameraPose}
{\texttt{camera\_name} (e.g. ``agentview'')}
{\texttt{pose}: Se3Pose}
{A named camera's pose in robot base frame.}

\skillcard{sim\_bridge.Init}
{\texttt{env\_name}; \texttt{config\_path} (optional); \texttt{camera\_names}; \texttt{task\_id}; \texttt{suite\_name}}
{\texttt{success}, \texttt{capabilities}: SimCapabilities, \texttt{error\_message}}
{Initialize a simulation environment by name.}

\skillcard{sim\_bridge.Reset}
{\texttt{seed}}
{\texttt{observation}: ObservationResponse}
{Reset the environment for a new episode.}

\skillcard{sim\_bridge.StepOnce}
{\texttt{gripper\_fraction} ($-1$ = no change)}
{\texttt{observation}, \texttt{reward}, \texttt{terminated}, \texttt{truncated}}
{Execute a single low-level MuJoCo timestep.}

\skillcard{sim\_bridge.MoveToJointsBlocking}
{\texttt{target\_joints}: JointState; \texttt{tolerance} (0.01); \texttt{max\_steps} (120)}
{\texttt{observation}, \texttt{reward}, \texttt{terminated}, \texttt{truncated}}
{Move to a joint configuration, blocking until converged.}

\skillcard{sim\_bridge.SetGripper}
{\texttt{fraction} (0 closed $\to$ 1 open); \texttt{arm\_id} (0)}
{(empty)}
{Set the gripper target fraction.}

\skillcard{sim\_bridge.GetState}
{(none)}
{\texttt{arm\_states}: ArmState[], \texttt{cameras}: CameraObservation[], \texttt{objects}: [name, pose] (ground truth)}
{Full sim state, including per-object ground-truth poses.}

\skillcard{sim\_bridge.CheckTaskCompletion}
{(none)}
{\texttt{reward}, \texttt{task\_completed}, \texttt{completion\_rate} (sub-goals satisfied)}
{Query simulator reward and task success.}

\skillcard{sim\_bridge.EnableVideoCapture}
{\texttt{enabled}; \texttt{clear}}
{(empty)}
{Start/stop video-frame capture.}

\skillcard{sim\_bridge.SaveVideo}
{\texttt{output\_path}; \texttt{clear}; \texttt{fps} (20)}
{\texttt{success}, \texttt{file\_path}, \texttt{num\_frames}}
{Encode captured frames to an MP4 and return its path.}
